\begin{theorem}[quantitative local maximum]\label{push:src:E.4.2}
If $|\tau-i|\leq3/20000$ and $\wp_\tau(z)=0$, then
$\mathcal N(z,\tau)\leq\sqrt2$, with equality precisely at
$\tau=i$. The explicit numerical second proof is retained in the compendium;
the global norm theorem gives the following short proof.
\end{theorem}
\Uses{\cref{ari:C.1.1,ari:C.8.15}}
\begin{proof}

For $\Lambda=\Z+\Z\tau$, the normalization identities are
\[
 \Delta_\Lambda=(2\pi)^{12}\eta_D(\tau)^{24},\qquad
 \sigma_\Lambda(z)=
 \frac{\vartheta_1(z,\tau)}{2\pi\eta_D(\tau)^3}
                 e^{\eta_1z^2/2}.
\]
The real-linear quasiperiod formula in \ref{ari:C.1.1} then gives
\[
 |\Delta_\Lambda|^{1/12}|\sigma_\Lambda(z)|
 e^{-\Re(z\eta(z))/2}
 =\left|\frac{\vartheta_1(z,\tau)}{\eta_D(\tau)}\right|
            e^{-\pi(\Im z)^2/\Im\tau}
 =\mathcal N(z,\tau).
\]
For $\wp_\tau(z)=0$, Theorem~\ref{ari:C.8.15} gives the
global maximum $\log\sqrt2$ of its logarithm, exactly on the
orbit of $i$. In the disc $|\tau-i|\le3/20000$ the only such
representative is $i$: a modular image with $c^2+d^2\ge2$
has imaginary part at most $1/2$, and the remaining images
are integral translates of $i$. Exponentiating proves the
stated inequality and its equality locus.

\end{proof}
